% =========================================================
% 08_determinants_inverse.tex
% Vectors and Linear Algebra
% =========================================================

% =========================================================
% SECTION DIVIDER
% =========================================================

\section{Determinants and Inverse Matrices}

% =========================================================
% DETERMINANT
% =========================================================

\begin{frame}{The Determinant}

The \textbf{determinant} is a scalar associated with a
square matrix.

For a $2\times2$ matrix,

\[
A=
\begin{bmatrix}
a&b\\
c&d
\end{bmatrix},
\]

the determinant is

\[
\boxed{
\det(A)=ad-bc
}
\]

\vspace{0.4cm}

The determinant provides information about:

\begin{itemize}
    \item whether a matrix is invertible,
    \item how a transformation changes area or volume,
    \item whether dimensions are collapsed by a transformation.
\end{itemize}

\end{frame}


% =========================================================
% DETERMINANT EXAMPLE
% =========================================================

\begin{frame}{Example: Determinant of a $2\times2$ Matrix}

Consider
\vspace{-0.1cm}
\[
A=
\begin{bmatrix}
3&2\\
1&4
\end{bmatrix}.
\]

Using

\[
\det(A)=ad-bc,
\]

we obtain

\[
\det(A)
=
(3)(4)-(2)(1)
=
12-2.
\]

Therefore,

\[
\boxed{
\det(A)=10
}
\]

Since the determinant is non-zero, $A$ is invertible.

\end{frame}


% =========================================================
% GEOMETRIC MEANING
% =========================================================

\begin{frame}{Geometric Meaning of the Determinant}

For a $2\times2$ transformation,

\[
|\det(A)|
\]

represents the factor by which areas are scaled.

\vspace{0.4cm}

\[
\begin{array}{c|c}
\text{Determinant} & \text{Geometric effect}\\
\hline
|\det(A)|>1 & \text{Area increases}\\[3pt]
0<|\det(A)|<1 & \text{Area decreases}\\[3pt]
|\det(A)|=1 & \text{Area is preserved}\\[3pt]
\det(A)=0 & \text{Area collapses to zero}
\end{array}
\]

\vspace{0.3cm}

The sign of the determinant also indicates whether
orientation is preserved or reversed.

\end{frame}


% =========================================================
% 3D DETERMINANT
% =========================================================

\begin{frame}{Determinant in Three Dimensions}

For a $3\times3$ matrix, the absolute value of the determinant
represents the factor by which volumes are scaled.

\[
\boxed{
|\det(A)|
=
\text{volume scaling factor}
}
\]

\vspace{0.4cm}

For example,

\[
\det(A)=3
\]

means that the transformation scales volume by a factor of $3$.

\vspace{0.4cm}

A negative determinant indicates a reversal of orientation.

\end{frame}


% =========================================================
% DETERMINANT PROPERTIES
% =========================================================

\begin{frame}{Properties of the Determinant}

For square matrices $A$ and $B$,

\[
\boxed{
\det(AB)=\det(A)\det(B)
}
\]

and

\[
\boxed{
\det(A^T)=\det(A).
}
\]

For the identity matrix,

\[
\boxed{
\det(I)=1.
}
\]

If $A$ is invertible,

\[
\boxed{
\det(A^{-1})
=
\frac{1}{\det(A)}.
}
\]

\end{frame}


% =========================================================
% SINGULAR MATRIX
% =========================================================

\begin{frame}{Singular and Non-Singular Matrices}

A square matrix $A$ is \textbf{singular} when

\[
\boxed{
\det(A)=0.
}
\]

A square matrix is \textbf{non-singular} when

\[
\boxed{
\det(A)\neq0.
}
\]

\vspace{0.4cm}

Geometrically, a singular transformation collapses the
space into a lower-dimensional space.

\vspace{0.3cm}

For example,

\[
\mathbb{R}^2
\longrightarrow
\text{a line}.
\]

\end{frame}


% =========================================================
% INVERTIBILITY
% =========================================================

\begin{frame}{Determinant and Invertibility}

For a square matrix $A$,

\[
\boxed{
A^{-1}\text{ exists}
\iff
\det(A)\neq0.
}
\]

Therefore,

\[
\boxed{
\det(A)=0
\iff
A\text{ is singular}.
}
\]

\vspace{0.4cm}

This gives an important connection:

\[
\boxed{
\text{Non-zero determinant}
\rightarrow
\text{Invertible matrix}
\rightarrow
\text{Inverse transformation}
}
\]

\end{frame}


% =========================================================
% INVERSE MATRIX
% =========================================================

\begin{frame}{Inverse Matrix}

The inverse of a square matrix $A$ is a matrix $A^{-1}$
such that

\[
\boxed{
AA^{-1}=A^{-1}A=I.
}
\]

\vspace{0.4cm}

The inverse matrix represents the transformation that
\textbf{undoes} the original transformation.

If

\[
\vect{y}=A\vect{x},
\]

then

\[
\boxed{
\vect{x}=A^{-1}\vect{y}.
}
\]

\vspace{0.3cm}

Thus, $A^{-1}$ maps the output back to the original input.

\end{frame}


% =========================================================
% 2x2 INVERSE FORMULA
% =========================================================

\begin{frame}{Inverse of a $2\times2$ Matrix}

For
\vspace{-0.1cm}
\[
A=
\begin{bmatrix}
a&b\\
c&d
\end{bmatrix},
\]

if
\vspace{-0.1cm}
\[
ad-bc\neq0,
\]

then
\vspace{-0.1cm}
\[
\boxed{
A^{-1}
=
\frac{1}{ad-bc}
\begin{bmatrix}
d&-b\\
-c&a
\end{bmatrix}.
}
\]

Notice that
\vspace{-0.1cm}
\[
ad-bc=\det(A).
\]

Therefore, the inverse exists precisely when

\[
\det(A)\neq0.
\]

\end{frame}


% =========================================================
% INVERSE EXAMPLE
% =========================================================

\begin{frame}{Example: Finding an Inverse}

Let
\vspace{-0.1cm}
\[
A=
\begin{bmatrix}
2&1\\
1&1
\end{bmatrix}.
\]

First,

\[
\det(A)
=
(2)(1)-(1)(1)
=
1.
\]

Therefore,

\[
A^{-1}
=
\begin{bmatrix}
1&-1\\
-1&2
\end{bmatrix}.
\]

Verification:

\[
AA^{-1}
=
\begin{bmatrix}
1&0\\
0&1
\end{bmatrix}
=
I.
\]

\end{frame}


% =========================================================
% INVERSE AND LINEAR SYSTEMS
% =========================================================

\begin{frame}{Inverse Matrix and Linear Systems}

Consider

\[
A\vect{x}=\vect{b}.
\]

If $A$ is invertible, multiply both sides by $A^{-1}$:

\[
A^{-1}A\vect{x}
=
A^{-1}\vect{b}.
\]

Since

\[
A^{-1}A=I,
\]

we obtain

\[
\boxed{
\vect{x}=A^{-1}\vect{b}.
}
\]

Thus, an invertible coefficient matrix gives a unique solution
for every $\vect{b}$.

\end{frame}


% =========================================================
% COFACTORS
% =========================================================

\begin{frame}{Minors and Cofactors}

For an element $a_{ij}$ of a matrix, its \textbf{minor}
$M_{ij}$ is the determinant obtained by deleting row $i$
and column $j$.

\vspace{0.4cm}

The corresponding cofactor is

\[
\boxed{
C_{ij}=(-1)^{i+j}M_{ij}.
}
\]

\vspace{0.4cm}

The signs follow the pattern

\[
\begin{bmatrix}
+&-&+\\
-&+&-\\
+&-&+
\end{bmatrix}.
\]

\vspace{0.3cm}

Cofactors can be used to calculate determinants of larger matrices.

\end{frame}


% =========================================================
% COFACTOR EXPANSION
% =========================================================

\begin{frame}{Cofactor Expansion}

For

\[
A=
\begin{bmatrix}
a_{11}&a_{12}&a_{13}\\
a_{21}&a_{22}&a_{23}\\
a_{31}&a_{32}&a_{33}
\end{bmatrix},
\]

expanding along the first row gives

\[
\boxed{
\det(A)
=
a_{11}C_{11}
+
a_{12}C_{12}
+
a_{13}C_{13}.
}
\]

\vspace{0.4cm}

Equivalently,

\[
\det(A)
=
a_{11}M_{11}
-a_{12}M_{12}
+a_{13}M_{13}.
\]

\vspace{0.3cm}

The same determinant can be expanded along any row or column.

\end{frame}


% =========================================================
% ADJUGATE
% =========================================================

\begin{frame}{Adjugate Matrix}

The matrix of cofactors is

\[
C=
\begin{bmatrix}
C_{11}&C_{12}&\cdots\\
C_{21}&C_{22}&\cdots\\
\vdots&\vdots&\ddots
\end{bmatrix}.
\]

The \textbf{adjugate} is the transpose of the cofactor matrix:

\[
\boxed{
\operatorname{adj}(A)=C^T.
}
\]

For an invertible matrix,

\[
\boxed{
A^{-1}
=
\frac{1}{\det(A)}
\operatorname{adj}(A).
}
\]

This provides a general formula for computing the inverse.

\end{frame}


% =========================================================
% LIMITATIONS
% =========================================================

\begin{frame}{Computing Inverses in Practice}

Although

\[
A^{-1}
=
\frac{1}{\det(A)}
\operatorname{adj}(A)
\]

is mathematically useful, it is not usually the preferred
method for large numerical problems.

\vspace{0.4cm}

For computational applications, methods such as

\begin{itemize}
    \item Gaussian elimination,
    \item LU factorization,
    \item QR factorization
\end{itemize}

are generally more practical.

\vspace{0.4cm}

The inverse formula is therefore particularly useful for
understanding the mathematics and for small matrices.

\end{frame}